\section{Auxiliary twists with prescribed local behavior}\label{sec:nonvanishing}

We need auxiliary twists with specified local behavior and nonzero
central value or derivative. For a fixed witness, no bound on the prime
count is needed. When the curve itself varies, we also need central-value
witnesses with a uniformly bounded number of prime factors. In both cases
their primes can be chosen beyond any previously specified finite set.

\begin{proposition}[Local witnesses]\label{nv:local-witness}
Fix an elliptic curve $A/\Q$ and a nonempty filter of fundamental
discriminants prime to $2N_A$, with specified real sign and unit
squareclasses at finitely many primes. If the root number in the filter
is $+1$, it contains $d$ of arbitrarily large absolute value with
$L(A^{(d)},1)\ne0$.
If the root number is $-1$, it contains such $d$ with
$L'(A^{(d)},1)\ne0$. Any additional finite set of primes may be excluded.
\end{proposition}

\begin{proposition}[Bounded-factor completion]\label{nv:rough}
There is an absolute integer $A_0$ with the following property.
In the positive-root-number case of Proposition~\ref{nv:local-witness},
one can require $\omega(d)\le A_0$ and require every prime divisor of
$d$ to exceed any fixed bound.
The size threshold for the resulting $d$ may depend on the curve, the
filter, and the excluded primes; $A_0$ does not.
\end{proposition}

The use of a nonnegative first moment and a lower linear sieve to obtain
bounded-factor twists goes back to Hoffstein--Luo
\cite[\S2]{HoffsteinLuo1997}. We retain an arbitrary fixed progression
and track an absolute factor-count bound without optimizing it.
The central moment follows the model of
\cite[Proposition~2]{RadziwillSoundararajan2015}, which has a sharper
divisor error; the derivative calculation adapts the first-moment
argument of \cite{Iwaniec1990}. The enlarged modulus and the absolute
exponents needed here are retained in the proof below.

A nonzero first moment alone supplies the derivative witness. For the
central-value completion, positivity allows us to sieve the first
moment, so we also need an error term for each prescribed squarefree
divisor of the twist parameter.

Write $f$ for the weight-two newform of $A$, normalized so that its
central point is $1/2$, and write $\lambda(n)$ for its normalized Hecke
eigenvalues. Thus
\[
L(s,f\otimes\chi)=\sum_{n\ge1}\lambda(n)\chi(n)n^{-s},
\qquad |\lambda(n)|\le d(n).
\]
Choose $H$ divisible by $8N_A$, a reduced residue class $a\bmod H$, and
$\sigma\in\{1,-1\}$ such that $\sigma a\equiv1\bmod4$.
The required twists have discriminant $\sigma n$ with
$n>0$ squarefree and $n\equiv a\bmod H$. A filter is a finite union of
such progressions, so it suffices to use one of them. Fix a nonnegative
smooth nonzero $\Phi$ supported in $(1,2)$.

\begin{lemma}[A fixed-progression moment]\label{nv:moment}
There are absolute constants $\eta>0$ and $B>0$ such that the
following assertions hold. For squarefree $b$ prime to $H$, put
\[
S_{j,b}(X)=
\sum_{\substack{n>0\ {\rm squarefree}\\ n\equiv a\ (H)\\b\mid n}}
 L^{(j)}(1/2,f\otimes\chi_{\sigma n})\Phi(n/X).
\]
If the common root number is $+1$, then
\[
S_{0,b}(X)=cX\prod_{p\mid b}\beta_p+
 O_{f,H,\Phi}(b^B X^{1-\eta}),
\qquad
c>0,\quad 0<\beta_p<1,\quad
\beta_p=\frac1p+O(p^{-2}).
\]
The last implied constant is absolute. If the common root number is
$-1$, then
\[
S_{1,1}(X)=cX\log X+c_1X+
 O_{f,H,\Phi}(X^{1-\eta}),\qquad c>0.
\]
\end{lemma}

\begin{proof}
We first control the square-divisor tail, then use Poisson summation to
separate a positive main term from a power-saving error. Tracking the
dependence on the imposed divisor $b$ will give the absolute exponents.
Let $G(z)=e^{z^2}$ and
\[
 V_j(y)=\frac1{2\pi i}\int_{(2)}
       G(z)\Gamma(1+z)y^{-z}\frac{dz}{z^{j+1}}.
\]
Writing $Q_A=\sqrt{N_A}/(2\pi)$, the symmetric approximate functional
equation in the required sign is
\begin{equation}\label{nv:afe}
L^{(j)}(1/2,f\otimes\chi_{\sigma n})
=2\sum_{u\ge1}\frac{\lambda(u)\chi_{\sigma n}(u)}{\sqrt u}
 V_j\!\left(\frac{u}{Q_A n}\right).
\end{equation}
For $j=1$, differentiating the completion contributes no central-value
term, because that value is zero.

\par\smallskip\noindent\emph{Square-divisor truncation.}
The quadratic large sieve
\cite[Theorem~1 and Corollary~1]{HeathBrown1995}, with the approximate functional equation,
gives for fixed $\delta>0$
\begin{equation}\label{nv:second}
 \sum_{\substack{0<t\le Z\ {\rm squarefree}\\
                 \sigma t\ {\rm fundamental}\\(t,2N_A)=1}}
 |L(1/2+\delta+i\tau,f\otimes\chi_{\sigma t})|^2
 \ll_{f,\rho,\delta} Z^{1+\rho}(1+|\tau|)^{C_\rho}.
\end{equation}
Here the polynomial exponent can be chosen independently of the level.
To see that square indices cause no extra power, write each coefficient
index as $ur^2$ with $u$ squarefree. Minkowski's inequality and the
quadratic large sieve bound the square root of the moment of a polynomial
of length $U$ by
\[
 (ZU)^{O(\rho')}
 \sum_{r\le\sqrt U}\frac{\sqrt{Z+U/r^2}}r
 \ll (ZU)^{O(\rho')}(\sqrt Z+\sqrt U).
\]
Coprimality restrictions at $r$ are discarded only after taking absolute
squares. To include both parts of the approximate functional equation,
work first with $t\asymp T$. Their lengths are
$U\ll_f T(1+|\tau|)^C$, with rapidly decreasing tails.
The main coefficients have the additional factor $n^{-\delta}$,
which improves the partial-summation bound. The dual coefficients
have $n^{\delta}$ and lose at most $U^{\delta}$, whereas their
functional-equation multiplier contributes $T^{-2\delta}$, up to
a fixed power of $1+|\tau|$. The resulting factor is at most
$T^{-\delta}(1+|\tau|)^{C_\delta}$ and causes no positive power
of $T$. Mellin inversion separates the smooth $t$-weights; the
gamma factors control the resulting integrals and tails.
Taking the auxiliary large-sieve exponents sufficiently small and
summing dyadic intervals proves Equation~\eqref{nv:second}.

Let $\mathcal A_j(n)$ denote the right side of
Equation~\eqref{nv:afe}, also for nonsquarefree $n\equiv a\bmod H$.
Expanding $\mu^2(n)$ and retaining divisors $d\le Y$ leaves an error
\begin{equation}\label{nv:tail}
 O_{f,H,\Phi,\rho}(X^{1+\rho}/Y).
\end{equation}
Indeed write $n=tv^2$, with $t$ squarefree. The omitted terms have $v>Y$,
and their multiplicity is at most the divisor count of $v$.
Let $L^{(v)}(s,f\otimes\chi_{\sigma t})$ be the primitive-twist
$L$-function with its Euler factors at primes dividing $v$ removed.
For $0<\delta_{\mathrm M}<\rho/8$, Mellin inversion gives
\[
 \mathcal A_j(tv^2)=\frac{2}{2\pi i}
 \int_{(\delta_{\mathrm M})}
 G(z)\Gamma(1+z)(Q_Atv^2)^z
 L^{(v)}(1/2+z,f\otimes\chi_{\sigma t})
 \frac{dz}{z^{j+1}}.
\]
Choose the large-sieve and divisor exponents smaller still.
On this line the removed Euler
factors cost an arbitrarily small fixed power of $v$. Cauchy's
inequality and Equation~\eqref{nv:second} give, for suitably small
$\rho'>0$,
\[
 \sum_{t\le2X/v^2}^{\!*}
 |L(1/2+\delta_{\mathrm M}+i\tau,f\otimes\chi_{\sigma t})|
 \ll_{f,\rho',\delta_{\mathrm M}}
 (X/v^2)^{1+\rho'}(1+|\tau|)^C,
\]
where the star retains the primitive discriminant conditions.
The gamma factors control the integral. The divisor count, the removed
Euler factors, and $X^{\delta_{\mathrm M}}$ are absorbed in $X^\rho$,
leaving $X^{1+\rho}\sum_{v>Y}v^{-2}\ll X^{1+\rho}/Y$.
The factor $z^{-2}$ changes only the constant.

\par\smallskip\noindent\emph{Poisson summation and the main term.}
Only $(d,H)=1$ occurs, because $n\equiv a\bmod H$ is a reduced
class. For $d\le Y$ put $J=[d^2,b]$ and
$c_J\equiv aJ^{-1}\bmod H$, which is therefore defined.
Write $u=u_1q$ with $u_1\mid H^\infty$ and $(q,H)=1$, and put
$\kappa(u_1)=\chi_{\sigma a}(u_1)$. Define
\[
 G_\nu(q)=\sum_{y\bmod q}\left(\frac yq\right)e(\nu y/q),
\quad
 W_j(z,R)=\int_0^\infty
 \Phi(x)V_j\!\left(\frac z{Q_Ax}\right)e(-Rx)\,dx,
\]
where $e(x)=e^{2\pi ix}$. Put
\[
 T_j(J)=\sum_{\substack{n>0,\ n\equiv a\ (H)\\J\mid n}}
             \mathcal A_j(n)\Phi(n/X).
\]
Poisson summation in the fixed progression gives the exact identity
\begin{align}\label{nv:poisson}
 T_j(J)
 &=\frac{2X}{HJ}
 \sum_{u_1\mid H^\infty}\frac{\lambda(u_1)\kappa(u_1)}{\sqrt{u_1}}
 \sum_{(q,H)=1}\frac{\lambda(q)}{q^{3/2}}
       \left(\frac{\sigma JH}{q}\right) \\
 &\hspace{14mm}\cdot
 \sum_{\nu\in\Z}G_\nu(q)e(\nu c_J\bar q/H)
 W_j\!\left(\frac{u_1q}{X},\frac{\nu X}{HJq}\right).
 \nonumber
\end{align}
Here $q\bar q\equiv1\bmod H$. The arithmetic factor follows directly by
applying the Chinese remainder theorem to
\[
\sum_{\substack{x\bmod Hq\\x\equiv c_J\ (H)}}
 \left(\frac xq\right)e(\nu x/(Hq))
=e(\nu c_J\bar q/H)\left(\frac Hq\right)G_\nu(q).
\]
When $(q,J)>1$, the additional Jacobi symbol is zero.

At zero frequency only square $q$ occur. Put
\[
 H_p(z)=1+(1-p^{-1})
     \sum_{r\ge1}\lambda(p^{2r})p^{-r(1+2z)},\qquad H_p=H_p(0),
\]
\[
 L_H(z)=\sum_{u_1\mid H^\infty}
             \lambda(u_1)\kappa(u_1)u_1^{-1/2-z},
 \qquad F_J(z)=\prod_{p\nmid HJ}H_p(z).
\]
The zero-frequency term is
\[
\frac{2X}{HJ}\frac1{2\pi i}\int_{(2)}
G(z)\Gamma(1+z)(Q_AX)^z\widehat\Phi_0(z)L_H(z)F_J(z)
\frac{dz}{z^{j+1}},
\quad
\widehat\Phi_0(z)=\int_0^\infty\Phi(x)x^z\,dx .
\]
For the Satake parameters $\alpha_p,\alpha_p^{-1}$, with
$|\alpha_p|=1$, we have
\[
 \sum_{r\ge0}\lambda(p^{2r})x^r
   =\frac{1+x}{(1-\alpha_p^2x)(1-\alpha_p^{-2}x)}.
\]
Thus $F_J$ factors through $L(1+2z,\operatorname{sym}^2f)$ and a
correction product convergent near $\Re z\ge-1/16$. Omitting a prime
dividing $J$ removes its correction factor and inserts an inverse
symmetric-square Euler polynomial; it creates no pole even when
$H_p(z)$ has zeros. The omitted factors cost at most $J^{B_1}$ for an
absolute $B_1$. Symmetric-square continuation, including the CM case
\cite[Theorem~9.3 and Remark~9.9]{GelbartJacquet1978}, and the
polynomial strip bounds from its functional equation therefore give an error
$O_{f,H,\Phi}(J^{B_1}X^{1-1/16})$ when shifting to $\Re z=-1/16$.

The residues have a positive leading coefficient. Indeed
\[
 F_p^{\rm ev}=\sum_{r\ge0}\lambda(p^{2r})p^{-r}
 =\frac{1+p^{-1}}{(1-\alpha_p^2/p)(1-\alpha_p^{-2}/p)}>0,\quad
 H_p=p^{-1}+(1-p^{-1})F_p^{\rm ev}>p^{-1}.
\]
The fixed Euler factors are positive at the center, and
$L(1,\operatorname{sym}^2f)>0$ by the Rankin--Selberg residue formula.
This also holds for CM forms. Summing the residue over the square
divisor gives, at $p\nmid Hb$, the factor $H_p-p^{-2}$, and at
$p\mid b$ the factor $p^{-1}-p^{-2}$. Consequently
\[
 \beta_p=\frac{p-1}{p^2H_p-1}
        =\frac1{1+pF_p^{\rm ev}}
        =p^{-1}+O(p^{-2}).
\]
For $j=1$, the double pole gives $cX\log X+c_1X$ with the same
positive leading coefficient. Extending the residue sum from $d\le Y$
to all $d$ costs $O(b^{B_2}X^{1+\rho}/Y)$; the $d^{-2+\rho}$ tail
also permits differentiation of that sum.

\par\smallskip\noindent\emph{The nonzero frequencies.}
It remains to obtain a power saving with polynomial dependence on $J$.
Put $\epsilon_q=1$ or $i$ according as $q\equiv1$ or $3\bmod4$
and $\widetilde G_\nu(q)=\epsilon_q^{-1}G_\nu(q)$. This is multiplicative
in $q$. Since $4\mid H$, expand
$\epsilon_q e(\nu c_J\bar q/H)$ in characters $\psi\bmod H$.
For $\nu\ne0$ the resulting Dirichlet series is
\[
 D_{\psi,J,\nu}(s)=
 \sum_{(q,H)=1}\frac{\lambda(q)}{q^s}
 \frac{\widetilde G_\nu(q)}{\sqrt q}
 \left(\frac{\sigma JH}{q}\right)\psi(q)
 =L(s,f\otimes\xi_{\rm prim})C_{\psi,J,\nu}(s),
\]
where $\xi(q)=\psi(q)(\sigma JH\nu/q)$.
At $p\nmid2HJ\nu$ its local factor is
$1+\lambda(p)\xi(p)p^{-s}$; dividing by the $\mathrm{GL}_2$ Euler
factor leaves a polynomial $1+O(p^{-2\Re s})$.
The bound at primes dividing $\nu$ has an absolute exponent.
Indeed, for an odd prime $p$, put $a=v_p(\nu)$. Directly summing
over the lifts of a residue modulo $p$ gives
\[
G_\nu(p^r)=
\begin{cases}
 \varphi(p^r),&r\text{ even},\ r\le a,\\
 -p^{r-1},&r\text{ even},\ r=a+1,\\
 \epsilon_p p^{r-1/2}
       \left(\dfrac{\nu/p^{r-1}}p\right),
       &r\text{ odd},\ r=a+1,\\
 0,&\text{otherwise}.
\end{cases}
\]
For even $r$ this is the Ramanujan sum; for odd $r$ the inner
sum of lifts vanishes unless $p^{r-1}\mid\nu$, after which
the primitive quadratic Gauss sum gives the displayed value.
In particular
$|\widetilde G_\nu(p^r)/p^{r/2}|\le p^{r/2}$, and the local
polynomial on $\Re s\ge3/4$ is bounded by
$\sum_{r\ge0}(r+1)p^{-r/4}\le C_0$ with $C_0$ absolute.
The inverse degree-two Euler polynomial is bounded by another
absolute constant $C_1$. Primes dividing $J$ simply remove the
local sum, and have the same bound after multiplication by that
inverse Euler polynomial. Thus all primes dividing $J\nu$ cost
at most $(C_0C_1)^{\omega(J\nu)}\le (2J|\nu|)^C$, for an
absolute $C$. The remaining correction product converges
absolutely on $\Re s\ge3/4$, since its local error is
$O(p^{-3/2})$. Hence the correction is holomorphic on
$\Re s>1/2$ and has the asserted polynomial bound with absolute exponent.
The conductor of $\xi_{\rm prim}$ is at most $4HJ|\nu|$, and the
twisted cuspidal $L$-function is entire. Thus, with absolute $C$,
\[
 |D_{\psi,J,\nu}(3/4+i\tau)|
 \ll_{f,H}(J|\nu|)^C(1+|\tau|)^C.
\]

Integration by parts in both weight variables gives, for fixed $r,s$
and any $K>0$,
\[
 |(z\partial_z)^r(R\partial_R)^sW_j(z,R)|
 \ll_{f,\Phi,r,s,K}(1+z)^{-K}(1+|R|)^{-K}(1+|\log z|).
\]
The logarithm is needed only for $j=1$ and is $O(\log X)$ on the
actual sums, where $z\ge X^{-1}$.
For an absolute small $\rho>0$ the sums can therefore be restricted to
\[
 u_1q\ll X^{1+\rho},\qquad
 0<|\nu|\ll_H JX^{2\rho},\qquad
 q\gg_H X^{1-\rho}/J,
\]
at a negligible cost with polynomial dependence on $J$.
A trivial bound $|G_\nu(q)|\le q$ gives a polynomial majorant for the
discarded tails. To identify the sum to be shifted, write
$D_{\psi,J,\nu}(s)=\sum_q\gamma(q)q^{-s}$. On a dyadic range
$q\asymp T$, take a fixed smooth cutoff $\vartheta$ and put
\[
 w_T(y)=\vartheta(y)
 W_j\!\left(\frac{u_1Ty}{X},\frac{\nu X}{HJTy}\right).
\]
Each dyadic $q$-sum after the character expansion is therefore
\[
 \sum_q\frac{\gamma(q)}q w_T(q/T)
 =\frac1{2\pi i}\int_{(1)}
       \widehat w_T(s)T^sD_{\psi,J,\nu}(1+s)\,ds,
\]
where $\widehat w_T(s)=\int_0^\infty w_T(y)y^{s-1}\,dy$.
The weight bounds permit a contour shift to $\Re s=-1/4$, giving
\[
 O_{f,H,\Phi}\bigl((J|\nu|)^C X^{C'\rho}T^{-1/4}\bigr).
\]
There is no residue. Summing over $\nu$, over the dyadic ranges,
and over $u_1$ (for which
$\sum_{u_1\mid H^\infty}|\lambda(u_1)|/\sqrt{u_1}<\infty$)
is bounded, up to a further $X^{O(\rho)}$, by
\[
 \frac XJ\,J^C(JX^{2\rho})^{C+1}
       (X^{1-\rho}/J)^{-1/4}
 =J^{2C+1/4}X^{3/4+O(\rho)}.
\]
This gives $O(J^{B_3}X^{3/4+C''\rho})$. Choose $\rho$ so that this is
$O(J^{B_3}X^{7/8})$.

\par\smallskip\noindent\emph{Choosing the truncation.}
Since $J\le bd^2$, summing the two contour errors for $d\le Y$,
and adjoining Equation~\eqref{nv:tail}, gives
\[
 O_{f,H,\Phi}\bigl(
 b^{B_4}Y^{2B_4+1}X^{1-\kappa}
 +b^{B_4}X^{1+\rho}/Y\bigr)
\]
for absolute $\kappa>0$ and $B_4$.
Take $Y=X^\delta$ with $(2B_4+1)\delta<\kappa/2$, and then
$\rho<\delta/4$. Both errors have the form in the statement,
with absolute positive saving. This proves the lemma.
\end{proof}

\begin{proof}[Proof of Propositions~\ref{nv:local-witness} and~\ref{nv:rough}]
For negative root number, Lemma~\ref{nv:moment} gives a nonzero sum of
derivatives in every sufficiently large dyadic interval. One derivative
is nonzero, and its central zero is already forced by the sign.

For positive root number, use the nonnegative sequence
\[
 a_n=\mu^2(n)\mathbf1_{n\equiv a\ (H)}
      L(1/2,f\otimes\chi_{\sigma n})\Phi(n/X).
\]
Its nonnegativity follows from the Waldspurger formula
\cite{Waldspurger1981}. The lemma gives a dimension-one sieve density
$\beta_p$ and an absolute exponent of distribution:
choose $\theta>0$ with $(B+1)\theta<\eta$ and put $D=X^\theta$.
The total remainder for squarefree $b\le D$ is
\[
 O_{f,H,\Phi}(D^{B+1}X^{1-\eta})=o(X/\log X).
\]
The density estimate also gives the interval product condition
\[
 \prod_{\substack{w\le p<z\\p\nmid H}}(1-\beta_p)^{-1}
 \le \frac{\log z}{\log w}
       \left(1+\frac{C_{f,H}}{\log w}\right)
 \qquad(2\le w<z).
\]
Indeed, comparison with $(1-p^{-1})^{-1}$ contributes a logarithmic
error $O(\sum_{p\ge w}p^{-2})$, and Mertens' product estimate gives
the displayed bound; finitely many small primes are absorbed in
$C_{f,H}$.  The same comparison gives
\[
 \prod_{\substack{p<z\\p\nmid H}}(1-\beta_p)\asymp_{f,H}(\log z)^{-1}.
\]
The lower linear sieve \cite[Theorem~11.13]{FriedlanderIwaniec2010},
with $z=D^{1/s}$ and any fixed $s>2$, therefore gives positive
remaining weight. Since $z=X^{\theta/s}$ tends to infinity, its
surviving parameters eventually avoid all primes up to any prescribed bound.
Every surviving parameter is prime to $H$ and has all other prime
factors at least $z$. Since $n<2X$, its number of prime factors is
bounded by any fixed integer greater than $2s/\theta$ for sufficiently
large $X$. This proves Proposition~\ref{nv:rough} and the positive-sign
part of Proposition~\ref{nv:local-witness}. Enlarging $H$ and choosing
a reduced CRT lift excludes any additional finite prime set.
\end{proof}

The constants in the size threshold may change when the curve or
the support is changed. The exponent $\theta$ and the factor-count
bound do not. In particular, a previously chosen twist can be
absorbed into the curve before this proposition is applied.
